% Part: first-order-logic
% Chapter: natural-deduction
% Section: rules-and-proofs

\documentclass[../../../include/open-logic-section]{subfiles}

\begin{document}

\iftag{FOL}
      {\olfileid{fol}{ntd}{rul}}
      {\olfileid{pl}{ntd}{rul}}

\olsection{Aturan lan \usetoken{P}{derivation}}

\begin{explain}
Sistem deduksi natural dimaksudake supaya cedhak karo penalaran ora
formal kang digunakake ing bukti matématika (mula rada ``natural'').
Bukti deduksi natural diwiwiti nganggo asumsi. Banjur aturan inferensi
ditrapake. Asumsi ``!!{discharged}'' dening aturan inferensi
\Intro{\lnot}, \Intro{\lif}, \iftag{FOL}{\Elim{\lor} lan
  \Elim{\lexists}}{lan \Elim{\lor}}, lan label asumsi kang
!!{discharged} diselehake ing sandhing inferensi supaya cetha.
\end{explain}

\begin{defn}[Asumsi]
\emph{Asumsi} yaiku sembarang !!{sentence}
ing posisi paling dhuwur ing sembarang cabang.
\end{defn}

!!^{derivation}s ing deduksi natural iku wit-wit !!{sentence}s
tartamtu: !!{sentence}s kang paling dhuwur iku asumsi, lan yen
!!a{sentence} mapan ing sangisore siji, loro, utawa telung ukara
liyane, ukara iku kudu ndherek kanthi bener miturut aturan
inferensi. Cathetan koreksi sumber OLPL-008: teks Inggris nyebut simpul
ing ndhuwur minangka sekuen, nanging definisi iki kanthi cetha
netepake kabeh simpul minangka ukara. !!{sentence}s ing sisih ndhuwur
inferensi diarani \emph{premis}, dene !!{sentence} ing sangisore
diarani \emph{kesimpulan} inferensi. Aturan-aturan iku awujud pasangan:
saben !!{operator} duwe aturan introduksi lan aturan eliminasi.
Aturan kasebut ngenalake !!a{operator} ing kesimpulan utawa ngilangi
!!a{operator} saka premis aturan. Sawetara aturan ngidini asumsi
kanthi jinis tartamtu kanggo \emph{!!{discharged}}. Kanggo nuduhake
asumsi endi kang !!{discharged} dening inferensi endi, kita uga menehi
label marang asumsi lan inferensi. Iki dituduhake kanthi nulis asumsi
minangka ``$\Discharge{!A}{n}$.''

% Only include this sentence is on of \land, lor, lif or lnot is defined.

\iftag{notprvNot,notprvAnd,notprvOr,notprvIf}{}%
	{Wis lumrah nimbang aturan kanggo kabeh !!{operator}s, yaiku $\land$, $\lor$, $\lif$, $\lnot$, lan $\lfalse$, sanajan sawetara operator kasebut ditetepake.}



\end{document}
